%%%%%%%%%%%%%%%%%%%%%%%%%%%%%%%%%%%%%%%%%%%%%%%%%%%%%%%%%%%%%%%
%                                                             %
% Cette version repart de 03_Draft\rolefirm_current-short.tex %
%                                                             %
%%%%%%%%%%%%%%%%%%%%%%%%%%%%%%%%%%%%%%%%%%%%%%%%%%%%%%%%%%%%%%%

% NB: Je supprime toutes les (grosses) parties qui étaient en commentaire dans la version précédente pour qu'on y voit plus clair.


Displaced workers experience a persistent decline in firm wage premia. This finding could reflect reallocation toward low-productivity firms but is also compatible with reemployment by productive firms with unfavorable wage policies despite being reasonably productive %This question is important in order to understand whether the reallocation process is productivity-enhancing or not.

In random-search job-ladder models with on-the-job search \citep[see][for a survey]{Lentz2024}, workers climb a ladder via outside offers and are pushed back down by displacement. The AKM firm wage premium identifies a firm-specific component of wages common to all workers of the firm, and is strongly positively correlated in the cross-section with productivity (Appendix Figure~\ref{fig:bscatter_firmfe_prod} and Table~\ref{tab:firm_fe_desc}; in line with \citealt{card2016bargaining} in Portugal). Structural displacement models in this class, such as \citet{krolikowski2017job} and \citet{Audoly_2024}, attribute the firm-side wage losses to reallocation toward less productive firms. Whether this productivity channel rationalizes the wage-premium decline we document is the empirical question we examine in what follows.


We test this hypothesis by looking at the productivity of firms in which displaced workers are re-employed: we estimate equation (\ref{eq:main_estimating_equation}), where $Y_{i(J)t}$ is a measure of productivity of firm $J$ where worker $i$ is employed in year $t$, computed as a three-year average centered on $t^D_i - 4$ (i.e.\ averaged over $t^D_i - 5$, $t^D_i - 4$, and $t^D_i - 3$). For each worker, firm characteristics are therefore measured over a fixed pre-displacement window relative to the worker's own event year. Accordingly, the event-study coefficients solely reflect the displacement-induced reallocation of workers toward firms with different \textit{initial} productivity rather than within-firm changes---similarly to how the analysis of the firm-specific wage premia is typically conducted in the literature. The corresponding coefficients are presented in Figure \ref{fig:ln_apl_ls} (thick dark lines) and in Table \ref{tab:evstudytable_body} (panel b).\footnote{A natural complementary question is how the productivity of laying-off firms themselves evolves around the event. Appendix Figure~\ref{fig:firm_es_productivity} reports a firm-level event-study showing that origin firms experience a temporary productivity decline around the layoff year and recover to slightly above their pre-layoff baseline by $h \geq +2$, consistent with the idea that layoffs allow firms to restore competitiveness.}


Figure \ref{fig:prod_res_within_btw} shows the impact of displacement on the average productivity (adjusted for capital intensity) of the subsequent employers.  %Firm-level productivity is measured over the 2001-2004 period and is held fixed, so that the event study coefficients solely reflect the displacement-induced reallocation of workers toward firms with different initial productivity.
Displaced workers reallocate on average toward significantly more productive firms than their counterfactual ($+0.14$ log points at $t^D+5$).\footnote{The same qualitative conclusion holds when using value-added per worker as an outcome (Table \ref{tab:evstudytable_body},column 2, panel b).} To gauge the magnitude of this effect, the employment-weighted interquartile range of productivity across firms is 0.53 -> \textcolor{red}{0.607} log points (Appendix Table~\ref{apptab:dispersion}). The effect on productivity at $t^D+5$ thus amounts to a quarter -> \textcolor{red}{22\%} of this range, equivalent to a move from the median of the productivity distribution to about its 63rd -> \textcolor{red}{60th} percentile. As expected given the average loss in firm wage premia, we further find that destination firms also feature a lower labor share of value added ($-0.07$ log points at $t^D+5$, Figure \ref{fig:ln_Lshare_within_btw}), in line with the notion that displaced workers reallocate towards higher-productivity but lower-paying firms.\footnote{This finding is consistent with the idea that market shares reallocation contributes to overall change in the labor share \citep{autor2020fall,aghion2022theory}.} The labor-share gap is 21\% of the employment-weighted interquartile range of the labor share across firms, equivalent to a move from the median to about the 40th -> \textcolor{red}{41st} percentile of the labor-share distribution. This inability to extract a larger share of value added is consistent with a loss in workers' negotiation power, which we explore in section \ref{sec:Nego}.

We assess the robustness of these findings in two ways. We first show that the same reallocation towards high-productivity, low-labor share, applies to the MLO sample (appendix Figure \ref{fig:prod_et_al_MLO}). Second, we compute the average of productivity and labor share over the full 2002--2015 period rather than over the three-year pre-displacement window centered on $t^D-4$, and show that the same qualitative pattern holds (appendix Figures \ref{fig:prod_et_al_Full_LE} and \ref{fig:prod_et_al_Full_MLO} for the LE and MLO samples respectively).

\label{par:accounting_consistency}A natural question is whether the firm-level findings reported above are quantitatively consistent with the $-4.2$ log-point AKM premium loss \textcolor{orange}{[$-3.7$]} documented in Section~\ref{sec:FEtheory}. The answer is not given by a mechanical adding-up exercise, because the firm-level outcomes (VA/$L$, labor share) are raw firm averages that embed firm-level technology together with workforce composition and other firm-level shocks, whereas $\psi_f$ is estimated net of worker fixed effects. We can nonetheless write the accounting identity $\Delta \ln \bar{w}_f = \Delta \psi_f + \Delta r_f$, where $\bar{w}_f$ is the firm's average wage and $r_f$ is a firm-level residual capturing workforce composition and other unobservables. At $t^D+5$, the labor share identity $\Delta \ln(\text{labor share}) = \Delta \ln \bar{w}_f - \Delta \ln(\text{VA}/L)$ implies $\Delta \ln \bar{w}_f \approx 0$, so that $\Delta \psi_f = -0.042$ is offset by $\Delta r_f = +0.042$. \textcolor{orange}{[R2: $\Delta\ln(\text{labor share}) = -0.074$ and $\Delta\ln(\text{VA}/L) = +0.066$ (Table 2, panel b, $t^D+5$) give $\Delta \ln \bar{w}_f \approx -0.008$, so $\Delta \psi_f = -0.037$ is offset by $\Delta r_f \approx +0.029$]} The drop in firm wage premia thus accounts for about $43\%$ of the $9.8$-log-point hourly wage scar \textcolor{orange}{[37\% of the 10.1-log-point scar: Table 2, cols. 5--7]}, with the residual term reflecting the difference in workforce composition and other unobservables between destination and origin firms. (\textcolor{red}{pw: numbers to review}) \textcolor{orange}{[see the orange notes in this paragraph]}

